\documentclass{article}
\pdfpagewidth=8.5in
\pdfpageheight=11in

\usepackage{WEDTreport}
% Use the postscript times font!
\usepackage{times}
\usepackage{soul}
\usepackage{url}
\usepackage{xcolor}
\usepackage{polski}
\usepackage[polish]{babel}
\usepackage[utf8]{inputenc}
\usepackage[T1]{fontenc}
\usepackage[utf8]{luainputenc}
\usepackage[hidelinks]{hyperref}
\usepackage[utf8]{inputenc}
\usepackage{caption}
\usepackage{indentfirst}
\usepackage{graphicx}
\usepackage{amsmath}
\usepackage{siunitx}
\usepackage{booktabs}
\usepackage{subfig}
\urlstyle{same}

\title{Wprowadzenie do eksploracji danych tekstowych w sieci WWW\\ Odpowiadanie na pytania}

\author{
Maria Konieczka, Alicja Poturała, Jakub Sikora
\affiliations
numery albumów: 283410, 283415, 283418 \\
\emails
maria.konieczka.stud@pw.edu.pl, alicja.poturala.stud@pw.edu.pl, jakub.sikora2.stud@pw.edu.pl
}

\newcommand{\todo}[1]{\textcolor{blue}{\textbf{TO DO:} #1}}

\begin{document}
\maketitle

\section{Wprowadzenie}
\label{sec:wprowadzenie}
Lorem ipsum dolor sit amet, consectetur adipiscing elit. In ultricies justo id elementum rutrum. Donec et diam turpis. Suspendisse mollis id enim a ultricies. Sed eleifend neque eu gravida elementum. Sed faucibus fermentum faucibus. Cras vitae vehicula orci, vitae convallis ante. Phasellus at odio ultrices, interdum lacus non, dignissim augue.

Quisque et mi sit amet ligula volutpat molestie eu in est. Sed dictum non augue eu ornare. Morbi tristique arcu ex, non dapibus lacus sollicitudin suscipit. Donec ante justo, vulputate vel lectus sed, lacinia pharetra libero. Nulla elementum erat augue, id porta lacus sagittis nec. Sed ultricies augue neque, a scelerisque lectus dictum ac. Sed fermentum faucibus sem vitae pretium. Morbi lobortis tortor vitae enim elementum aliquam. Aliquam mollis lacus eget ipsum dignissim, eget cursus tortor dapibus. Donec feugiat nunc ac placerat rhoncus. Vestibulum at tempus purus. Cras et pharetra orci, quis tempus augue.

Quisque non tortor metus. Vestibulum sem dui, tempor eu placerat non, pharetra eget metus. Vestibulum venenatis, est in aliquet placerat, mi tellus fermentum magna, ut congue ligula lectus placerat elit. Nullam id sollicitudin nisi. Suspendisse sodales hendrerit arcu posuere semper. Sed vel tortor accumsan, ullamcorper tortor eget, tincidunt nisl. Suspendisse auctor odio orci, pellentesque varius lectus gravida non. Aliquam in urna sed dolor sollicitudin euismod congue vitae risus. Ut malesuada massa massa. Suspendisse mollis et libero non eleifend. Duis at lobortis ante.

\begin{figure}[h]
    \centering
    \includegraphics[width=\columnwidth]{./figures/example.png}
    \caption{Morbi in erat tristique, ullamcorper felis eget, auctor risus}
    \label{fig:example-figure}
\end{figure}

Nulla ullamcorper nunc et dui sagittis lobortis. Aliquam erat volutpat. Proin congue purus at sapien aliquam, vel blandit enim fringilla. Curabitur facilisis turpis at ullamcorper vestibulum. Morbi in erat tristique, ullamcorper felis eget, auctor risus. Sed augue turpis, auctor cursus leo in, lacinia finibus enim. Sed ac tortor ac felis dignissim ullamcorper vel ut neque. Duis at arcu tortor. Nunc eu tincidunt lacus. Donec vitae nibh arcu. Phasellus gravida rhoncus diam viverra vulputate.

Aenean odio erat, rhoncus in pharetra sit amet, volutpat ac metus. Nam interdum quam nibh, vel semper lectus luctus sit amet. Vivamus eu dui dolor. Quisque euismod magna in dignissim faucibus. Curabitur vulputate ex id leo posuere, sed finibus quam ornare. Etiam eget imperdiet metus. Morbi odio est, iaculis pellentesque cursus et, interdum id erat. Aliquam ultricies metus risus, eu rutrum mauris fringilla sed. Duis malesuada nisi ex, molestie tristique metus tempor et. Curabitur vel est ac libero bibendum euismod. Maecenas vitae vulputate nisl, dignissim vehicula ipsum. Proin lobortis luctus augue, vel elementum erat sagittis varius. Mauris elementum, turpis vel fermentum vestibulum, ex turpis mattis eros, ut vehicula lorem nunc blandit enim. Suspendisse eget vestibulum mauris. Aenean fermentum tellus sit amet nulla pellentesque fringilla.

Curabitur ut mauris dignissim, fermentum lectus et, lacinia ligula. Cras at turpis nibh. Suspendisse interdum ipsum sed purus suscipit, ac auctor metus pretium. Vestibulum eu nisl ut magna pellentesque eleifend quis sed tellus. Nulla viverra magna non massa malesuada posuere. Cras sollicitudin fringilla lectus eget faucibus. Nam ac sapien consequat, consequat metus vel, feugiat ipsum. Interdum et malesuada fames ac ante ipsum primis in faucibus. Nulla ornare lectus augue, eget dictum nulla fringilla non. Aliquam at mauris eget sapien ornare aliquam~\cite{book-without-title}.

\bibliographystyle{abbrv}
\bibliography{bibliography}

\end{document}